\documentclass{article}
\usepackage[T1]{fontenc}
\usepackage{lmodern}
\usepackage{graphicx}
\usepackage{amsmath,amssymb}
\usepackage[a4paper,margin=2cm]{geometry}
\usepackage[absolute,overlay]{textpos}
\setlength{\TPHorizModule}{1mm}
\setlength{\TPVertModule}{1mm}
\usepackage[indonesian]{babel}
\usepackage{hyperref}
\setlength{\emergencystretch}{2em}
% unit: Halaman 41

\begin{document}

% === Halaman 41 (python/native) ===

8.  Sistem kriptografi (cryptosystem) 


    Sistem kriptografi adalah quintuple (E, D, M, K, C)

• M  adalah himpunan plainteks

• K  adalah himpunan kunci

• C  adalah himpunan cipherteks

• E  adalah himpunan fungsi enkripsi E: M  K → C

• D  adalah himpunan fungsi dekripsi D: C  K → M

41

\end{document}
